\renewcommand{\theposttitle}{2 Types of Burnout}
\renewcommand{\thepostdate}{2023-10-10}
\vspace*{\fill}
\begin{center}
\includegraphics[width=\textwidth, height=0.55\textheight, keepaspectratio]{C:/Users/James/Documents/GitHub/VisualIdeasBook/data/images/d3c342fc707ffcfa87e63aef80783852.png}
\end{center}
\vspace{2mm}
\begin{center}
\begin{minipage}{0.92\textwidth}
\small
Burning out can happen at 2 intensities:

Most people understand \textbf{“fast burnout”}—when you push yourself beyond capacity towards a goal, sacrificing your physical and mental health.

But equally dangerous is \textbf{“slow burnout”}—when you feel trapped in a situation and your motivation to do anything slowly suffocates. 

We often suffer silently from a misplaced belief that it's all part of "enduring the grind."

To dive deeper, my friend Roberto Ferraro shared this link:

\textbf{Harvard Business Review: }\\
"Is Your Burnout From Too Much Work or Too Little Impact?"
\end{minipage}
\end{center}
\vspace*{\fill}